\documentclass[12pt, a4paper]{article}
\usepackage[utf8]{inputenc}
\usepackage[spanish]{babel}
\usepackage{amsmath}
\usepackage{amssymb}
\usepackage{booktabs}
\usepackage{geometry}
\usepackage{xcolor}
\usepackage{graphicx}

\geometry{margin=2.5cm}

\title{\textbf{Prueba Calc}\\[0.5em]
\large LaTeX + celdas Python con \texttt{\textbackslash py\{\}}}
\author{}
\date{\today}

\begin{document}
\maketitle
\tableofcontents
\newpage

% ===========================================================
\section{Cálculos básicos}
% ===========================================================

%#python
%import math
%
%# Variables básicas
%radio       = 5.0
%altura      = 10.0
%gravedad    = 9.81   # m/s²
%masa        = 70.0   # kg
%
%# Cálculos
%area_circulo    = math.pi * radio**2
%perimetro       = 2 * math.pi * radio
%vol_cilindro    = area_circulo * altura
%peso            = masa * gravedad
%
%print(f"Área círculo  : {area_circulo:.4f} m²")
%print(f"Volumen cil.  : {vol_cilindro:.4f} m³")
%print(f"Peso          : {peso:.2f} N")
%#end

Para un círculo de radio $r = 5$ m:

\[
  A = \pi r^2 = 78.5398 \text{ m}^2
  \qquad
  P = 2\pi r = 31.4159 \text{ m}
\]

El volumen del cilindro de altura $h = 10$ m es:

\[
  V = A \cdot h = 78.5398 \times 10
    = 785.3982 \text{ m}^3
\]

Una persona de 70 kg pesa $686.7$ N
($0.687$ kN). También puedes operar directamente
en el documento: $2A - \tfrac{3}{2} = 155.579632679$.


% ===========================================================
\section{Secuencias y tablas}
% ===========================================================

%#python
%# Tabla de potencias (devuelve filas LaTeX como cadena)
%filas_potencias = "\n    ".join(
%    f"{n} & {n**2} & {n**3} & {n**4} \\\\"
%    for n in range(1, 8)
%)
%
%# Números de Fibonacci
%def fib(n):
%    a, b = 0, 1
%    seq = []
%    for _ in range(n):
%        seq.append(a)
%        a, b = b, a + b
%    return seq
%
%fib10 = fib(10)
%fib_str = ", ".join(str(x) for x in fib10)
%fib_sum = sum(fib10)
%
%print("Fibonacci (10):", fib_str)
%print("Suma:", fib_sum)
%#end

\subsection{Potencias}

\begin{table}[h]
  \centering
  \caption{Potencias de $n$ calculadas en Python}
  \begin{tabular}{c c c c}
    \toprule
    $n$ & $n^2$ & $n^3$ & $n^4$ \\
    \midrule
    1 & 1 & 1 & 1 \\
    2 & 4 & 8 & 16 \\
    3 & 9 & 27 & 81 \\
    4 & 16 & 64 & 256 \\
    5 & 25 & 125 & 625 \\
    6 & 36 & 216 & 1296 \\
    7 & 49 & 343 & 2401 \\
    \bottomrule
  \end{tabular}
\end{table}

\subsection{Sucesión de Fibonacci}

Los primeros 10 términos de la sucesión de Fibonacci son:

\[
  0, 1, 1, 2, 3, 5, 8, 13, 21, 34
\]

Su suma es $88$.


% ===========================================================
\section{Trigonometría}
% ===========================================================

%#python
%import math
%
%angulos_deg = [0, 30, 45, 60, 90]
%
%# Tabla trigonométrica (notación matemática: ^{\circ} para grados)
%filas_trig = ""
%for deg in angulos_deg:
%    rad = math.radians(deg)
%    s   = round(math.sin(rad), 4)
%    c   = round(math.cos(rad), 4)
%    t   = r"\infty" if deg == 90 else str(round(math.tan(rad), 4))
%    filas_trig += f"    ${deg}^{{\\circ}}$ & ${s}$ & ${c}$ & ${t}$ \\\\\n"
%
%# Identidad de Euler evaluada en pi
%euler_re = round(math.cos(math.pi), 6)   # = -1
%euler_im = round(math.sin(math.pi), 6)   # ≈ 0
%
%print("Tabla trig. generada.")
%print(f"e^(iπ) = {euler_re} + {euler_im}i")
%#end

\subsection{Tabla trigonométrica}

\begin{table}[h]
  \centering
  \caption{Valores de seno, coseno y tangente}
  \begin{tabular}{c c c c}
    \toprule
    Ángulo & $\sin$ & $\cos$ & $\tan$ \\
    \midrule
    $0^{\circ}$ & $0.0$ & $1.0$ & $0.0$ \\
    $30^{\circ}$ & $0.5$ & $0.866$ & $0.5774$ \\
    $45^{\circ}$ & $0.7071$ & $0.7071$ & $1.0$ \\
    $60^{\circ}$ & $0.866$ & $0.5$ & $1.7321$ \\
    $90^{\circ}$ & $1.0$ & $0.0$ & $\infty$ \\

    \bottomrule
  \end{tabular}
\end{table}

\subsection{Identidad de Euler}

La famosa identidad $e^{i\pi} + 1 = 0$ verificada numéricamente:

\[
  e^{i\pi} = \cos(\pi) + i\sin(\pi)
           = -1 + 0\,i
\]


% ===========================================================
\section{Estadística}
% ===========================================================

%#python
%import math
%
%datos = [12.3, 15.7, 9.8, 22.1, 18.4, 11.2, 25.6, 14.9, 17.3, 20.0]
%n     = len(datos)
%media = sum(datos) / n
%varianza = sum((x - media)**2 for x in datos) / n
%desv_std = math.sqrt(varianza)
%minimo   = min(datos)
%maximo   = max(datos)
%rango    = maximo - minimo
%datos_ord = sorted(datos)
%mediana = (datos_ord[n//2 - 1] + datos_ord[n//2]) / 2
%
%datos_str = ", ".join(f"{x}" for x in datos)
%
%print(f"N={n}  media={media:.3f}  σ={desv_std:.3f}  mediana={mediana:.1f}")
%#end

Para el conjunto de datos $\{\ 12.3, 15.7, 9.8, 22.1, 18.4, 11.2, 25.6, 14.9, 17.3, 20.0\ \}$ ($n = 10$):

\[
  \bar{x} = 16.730
  \qquad
  \sigma  = 4.732
  \qquad
  \tilde{x} = 16.5
\]

Rango: $[9.8,\ 25.6]$, amplitud $= 15.8$.


% ===========================================================
\section{Expresiones directas en el texto}
% ===========================================================

También puedes escribir expresiones Python directamente en el texto
sin declarar nada en una celda:

\begin{itemize}
  \item $\sqrt{2} = 1.414214$
  \item $\ln(e^3) = 3$
  \item $100! \mod 97 = 0$
  \item Longitud de la lista anterior: 10 elementos
\end{itemize}

\end{document}
